\section{Prompt Templates Used to Query S-LLMs}
\label{app:prompt_templates}

\noindent The templates below are fixed across all evaluated S-LLMs to ensure reproducibility; only the bracketed fields (\texttt{\{first\_prompt\}}, \texttt{\{code\}}, \texttt{\{prompt\}}) are instantiated per task/turn.

\begin{PromptBlock}{Turn 1 Wrapper Prompt (DT construction from a vague request)}
You are a PyChrono expert tasked with generating a simulation script based on the following instructions. Make sure to:
(*@\textbf{1.}@*) Initialize the PyChrono environment and core components.\\
(*@\textbf{2.}@*) Add the required physical systems and objects as specified.\\
(*@\textbf{3.}@*) Set necessary default parameters such as positions, forces, and interactions when not explicitly provided.\\
(*@\textbf{4.}@*) Include a simulation loop (and visualization if requested).\\
(*@\textbf{5.}@*) Output only the final runnable Python script (no explanations, no markdown).

Instructions:
(*@\texttt{\{first\_prompt\}}@*)
\end{PromptBlock}

\begin{PromptBlock}{Turns 2--3 Wrapper Prompt (code editing with bug fixing and sharp modifications)}
You are a PyChrono expert tasked with generating a simulation script based on the following instructions and a given PyChrono script, which may contain errors.

Your task has two parts:
(*@\textbf{(i)}@*) Identify and correct any issues that would prevent correct execution (including syntax errors, logical errors, incorrect method names, and parameter issues).\\
(*@\textbf{(ii)}@*) Modify the script to meet the new requirements.

Preserve existing functionality unless it directly conflicts with the new requirements.

IMPORTANT: Output only the final corrected and modified runnable Python script (no explanations, no markdown).

Here is the PyChrono code you need to modify:
(*@\texttt{\{code\}}@*)

Please modify the given code based on the following instructions:
(*@\texttt{\{prompt\}}@*)
\end{PromptBlock}


\section{J-LLM Prompt Templates}
\label{app:jllm_prompts}

\noindent The following templates are used to query the rubric-based judge model (J-LLM). 
Bracketed fields (\texttt{\{code\}}, \texttt{\{reference\_code\}}, \texttt{\{api\_documentation\_link\}}) are instantiated per task/turn.

% ---------------------------
% J-LLM_Ref_Doc
% ---------------------------
\begin{PromptBlock}{J-LLM\_Ref\_Doc (Reference + API Documentation)}
You are a PyChrono expert tasked with evaluating a simulation script by comparing it against a reference script generated by experts. Your evaluation should consider both the accuracy of the script compared to the reference and adherence to best practices as outlined in the PyChrono API documentation.

Here is the PyChrono code you need to evaluate:
[The Start of Assistant's Answer]
(*@\texttt{\{code\}}@*)
[The End of Assistant's Answer]

Here is the expert-generated reference code:
[The Start of Reference Answer]
(*@\texttt{\{reference\_code\}}@*)
[The End of Reference Answer]

Use the following evaluation criteria and point deduction guidelines:

(*@\textbf{1. Completeness (40 points total)}@*)
- Compare the provided code to the reference script. Deduct (*@\textbf{15 points}@*) for each missing essential component (e.g., system initialization, body creation, visualization) that is present in the reference script.
- Deduct (*@\textbf{10 points}@*) if a component is present but lacks important details or is incorrectly configured compared to the reference.
- Deduct (*@\textbf{5 points}@*) for minor omissions or slight deviations from the reference.

(*@\textbf{2. Correctness (30 points total)}@*)
- Compare the code to the reference script. Deduct (*@\textbf{15 points}@*) for each incorrect use of a PyChrono API that could lead to a significant change in simulation behavior.
- Deduct (*@\textbf{10 points}@*) for logical errors in the code, such as incorrect joint initialization or wrong setting of body properties, especially if the reference script does it correctly.
- Deduct (*@\textbf{5 points}@*) for minor inaccuracies or unnecessary API calls that deviate from the reference script.

(*@\textbf{3. Code Quality (10 points total)}@*)
- Evaluate the readability, structure, and documentation of the code against the reference script. Deduct (*@\textbf{5--10 points}@*) for poor readability, structure, or lack of meaningful variable names and formatting.
- Deduct (*@\textbf{5 points}@*) for insufficient comments or failure to follow documentation best practices, especially if the reference script provides better documentation.

(*@\textbf{4. Efficiency (10 points total)}@*)
- Evaluate the efficiency of the code compared to the reference script. Deduct (*@\textbf{5 points}@*) for each instance of unnecessary calculations, redundant code, or inefficient use of APIs that is optimized in the reference script.
- Deduct (*@\textbf{3 points}@*) for missing obvious optimization opportunities that the reference script implements.

(*@\textbf{5. Error Handling and Robustness (5 points total)}@*)
- Assess the error handling and robustness of the code. Deduct (*@\textbf{5 points}@*) for lack of basic error handling or failure to account for common issues that the reference script handles.
- Deduct (*@\textbf{3 points}@*) for inadequate handling of edge cases compared to the reference script.

(*@\textbf{6. Use of Visualization Tools (5 points total)}@*)
- Compare the use of visualization tools in the provided code to the reference script. Deduct (*@\textbf{3--5 points}@*) for incorrect or inadequate visualization setup as per the reference script.
- Deduct (*@\textbf{2 points}@*) for minor visualization issues, such as suboptimal lighting or incomplete setup of visual elements, compared to the reference.

Avoid position biases and ensure that the order in which the responses are presented does not influence your decision. Do not allow the length of the responses to influence your evaluation. Do not favor certain names of the assistants. Be as objective as possible.

Reference the PyChrono API documentation provided here: (*@\texttt{\{api\_documentation\_link\}}@*)

After providing your explanation, output the final score using the following format: (*@\textbf{[[x]]}@*) where x is the score assigned to the assistant's answer.

Provide the evaluated score and a brief explanation of the deductions below:
\end{PromptBlock}

% ---------------------------
% J-LLM_Ref
% ---------------------------
\begin{PromptBlock}{J-LLM\_Ref (Reference Only)}
You are a PyChrono expert tasked with evaluating a simulation script by comparing it against a reference script generated by experts.

Here is the PyChrono code you need to evaluate:
[The Start of Assistant's Answer]
(*@\texttt{\{code\}}@*)
[The End of Assistant's Answer]

Here is the expert-generated reference code:
[The Start of Reference Answer]
(*@\texttt{\{reference\_code\}}@*)
[The End of Reference Answer]

Use the following evaluation criteria and point deduction guidelines:

(*@\textbf{1. Completeness (40 points total)}@*)
- Compare the provided code to the reference script. Deduct (*@\textbf{15 points}@*) for each missing essential component (e.g., system initialization, body creation, visualization) that is present in the reference script.
- Deduct (*@\textbf{10 points}@*) if a component is present but lacks important details or is incorrectly configured compared to the reference.
- Deduct (*@\textbf{5 points}@*) for minor omissions or slight deviations from the reference script.

(*@\textbf{2. Correctness (30 points total)}@*)
- Compare the code to the reference script. Deduct (*@\textbf{15 points}@*) for each incorrect use of a PyChrono API that could lead to a significant change in simulation behavior.
- Deduct (*@\textbf{10 points}@*) for logical errors in the code, such as incorrect joint initialization or wrong setting of body properties, especially if the reference script does it correctly.
- Deduct (*@\textbf{5 points}@*) for minor inaccuracies or unnecessary API calls that deviate from the reference script.

(*@\textbf{3. Code Quality (10 points total)}@*)
- Evaluate the readability, structure, and documentation of the code against the reference script. Deduct (*@\textbf{5--10 points}@*) for poor readability, structure, or lack of meaningful variable names and formatting.
- Deduct (*@\textbf{5 points}@*) for insufficient comments or failure to follow documentation best practices, especially if the reference script provides better documentation.

(*@\textbf{4. Efficiency (10 points total)}@*)
- Evaluate the efficiency of the code compared to the reference script. Deduct (*@\textbf{5 points}@*) for each instance of unnecessary calculations, redundant code, or inefficient use of APIs that is optimized in the reference script.
- Deduct (*@\textbf{3 points}@*) for missing obvious optimization opportunities that the reference script implements.

(*@\textbf{5. Error Handling and Robustness (5 points total)}@*)
- Assess the error handling and robustness of the code. Deduct (*@\textbf{5 points}@*) for lack of basic error handling or failure to account for common issues that the reference script handles.
- Deduct (*@\textbf{3 points}@*) for inadequate handling of edge cases compared to the reference script.

(*@\textbf{6. Use of Visualization Tools (5 points total)}@*)
- Compare the use of visualization tools in the provided code to the reference script. Deduct (*@\textbf{3--5 points}@*) for incorrect or inadequate visualization setup as per the reference script.
- Deduct (*@\textbf{2 points}@*) for minor visualization issues, such as suboptimal lighting or incomplete setup of visual elements, compared to the reference.

Avoid position biases and ensure that the order in which the responses are presented does not influence your decision. Do not allow the length of the responses to influence your evaluation. Do not favor certain names of the assistants. Be as objective as possible.

After providing your explanation, output the final score using the following format: (*@\textbf{[[x]]}@*) where x is the score assigned to the assistant's answer.

Provide the evaluated score and a brief explanation of the deductions below:
\end{PromptBlock}

% ---------------------------
% J-LLM_Doc
% ---------------------------
\begin{PromptBlock}{J-LLM\_Doc (API Documentation Only)}
You are a PyChrono expert tasked with evaluating a simulation script by comparing it against the PyChrono API documentation. While the API documentation provides guidelines, it may not cover all aspects due to length constraints. Therefore, your evaluation should also be based on your knowledge of best practices in Python coding and general simulation principles.

Here is the PyChrono code you need to evaluate:
[The Start of Assistant's Answer]
(*@\texttt{\{code\}}@*)
[The End of Assistant's Answer]

Use the following evaluation criteria and point deduction guidelines:

(*@\textbf{1. Completeness (40 points total)}@*)
- Deduct (*@\textbf{15 points}@*) for each missing essential component (e.g., system initialization, body creation, visualization) as outlined in the PyChrono API documentation or generally expected in a simulation setup.
- Deduct (*@\textbf{10 points}@*) if a component is present but lacks important details or is incorrectly configured according to the API documentation or general simulation best practices.
- Deduct (*@\textbf{5 points}@*) for minor omissions or slight deviations from best practices mentioned in the API documentation or common Python coding practices.

(*@\textbf{2. Correctness (30 points total)}@*)
- Deduct (*@\textbf{15 points}@*) for each incorrect use of a PyChrono API that could lead to a significant change in simulation behavior, as indicated by the documentation or your expert knowledge.
- Deduct (*@\textbf{10 points}@*) for logical errors in the code, such as incorrect joint initialization or wrong setting of body properties, based on the API documentation or standard simulation principles.
- Deduct (*@\textbf{5 points}@*) for minor inaccuracies or unnecessary API calls that deviate from the API guidelines or standard coding practices.

(*@\textbf{3. Code Quality (10 points total)}@*)
- Evaluate the readability, structure, and documentation of the code. Deduct (*@\textbf{5--10 points}@*) for poor readability, structure, or lack of meaningful variable names and formatting, based on your Python expertise.
- Deduct (*@\textbf{5 points}@*) for insufficient comments or failure to follow documentation best practices, whether outlined in the API documentation or based on general coding standards.

(*@\textbf{4. Efficiency (10 points total)}@*)
- Deduct (*@\textbf{5 points}@*) for each instance of unnecessary calculations, redundant code, or inefficient use of APIs that could be optimized according to the API documentation or your understanding of efficient coding practices.
- Deduct (*@\textbf{3 points}@*) for missing obvious optimization opportunities as suggested by the API documentation or standard programming practices.

(*@\textbf{5. Error Handling and Robustness (5 points total)}@*)
- Deduct (*@\textbf{5 points}@*) for lack of basic error handling or failure to account for common issues, as recommended by the API documentation or best practices in Python coding.
- Deduct (*@\textbf{3 points}@*) for inadequate handling of edge cases, considering both the API documentation and typical robustness requirements in coding.

(*@\textbf{6. Use of Visualization Tools (5 points total)}@*)
- Deduct (*@\textbf{3--5 points}@*) for incorrect or inadequate visualization setup according to the API documentation or general expectations for visualizing simulations.
- Deduct (*@\textbf{2 points}@*) for minor visualization issues, such as suboptimal lighting or incomplete setup of visual elements, based on both the API documentation and your understanding of effective simulation visualization.

Avoid position biases and ensure that the order in which the responses are presented does not influence your decision. Do not allow the length of the responses to influence your evaluation. Do not favor certain names of the assistants. Be as objective as possible.

Reference the PyChrono API documentation provided here: (*@\texttt{\{api\_documentation\_link\}}@*)

After providing your explanation, output the final score using the following format: (*@\textbf{[[x]]}@*) where x is the score assigned to the assistant's answer.

Provide the evaluated score and a brief explanation of the deductions below:
\end{PromptBlock}


\section{Curiosity Rover Example: Ground-Truth (Expert Reference) Code}
\label{app:curiosity_gt}

\noindent The following scripts are expert-authored reference implementations for the Curiosity rover example used in paragraph~\ref{app:curiosity_example}. They correspond to Turns~1--3 respectively.

\begin{CodeBlock}{Curiosity (RBT) --- Turn 1 Ground Truth (Expert Reference)}
import os
import math
import numpy as np
import pychrono as chrono
import pychrono.robot as robot
from pychrono import irrlicht as chronoirr

# Create Chrono system
system = chrono.ChSystemNSC()
system.SetCollisionSystemType(
    chrono.ChCollisionSystem.Type_BULLET)
system.SetGravitationalAcceleration(chrono.ChVector3d(0, 0, -9.81))
chrono.ChCollisionModel.SetDefaultSuggestedEnvelope(
    0.0025)
chrono.ChCollisionModel.SetDefaultSuggestedMargin(0.0025)

# Create ground body with contact material and add it to the system
ground_mat = chrono.ChContactMaterialNSC()
ground = chrono.ChBodyEasyBox(20, 20, 1, 1000, True, True, ground_mat)
ground.SetPos(chrono.ChVector3d(0, 0, -0.5))  # Position the ground slightly below the origin
ground.SetFixed(True)  # Fix the ground in place
ground.GetVisualShape(0).SetTexture(
    chrono.GetChronoDataFile("textures/concrete.jpg"))
system.Add(ground)

# Create Curiosity rover and add it to the system
rover = robot.Curiosity(system)

# Create driver for rover
driver = robot.CuriosityDCMotorControl()
rover.SetDriver(driver)

# Initialize rover position and orientation
init_pos = chrono.ChVector3d(0, 0.2, 0)
init_rot = chrono.ChQuaterniond(1, 0, 0, 0)
rover.Initialize(chrono.ChFramed(init_pos, init_rot))

# Create the Irrlicht visualization
vis = chronoirr.ChVisualSystemIrrlicht()
vis.AttachSystem(system)
vis.SetCameraVertical(chrono.CameraVerticalDir_Z)
vis.SetWindowSize(1280, 720)
vis.SetWindowTitle("Curiosity rover - Rigid terrain")
vis.Initialize()
vis.AddLogo(chrono.GetChronoDataFile(
    "logo_pychrono_alpha.png"))
vis.AddSkyBox()
vis.AddCamera(chrono.ChVector3d(0, 3, 3), chrono.ChVector3d(0, 0, 0))
vis.AddTypicalLights()
vis.AddLightWithShadow(chrono.ChVector3d(1.5, -2.5, 5.5), chrono.ChVector3d(0, 0, 0), 3, 4, 10, 40, 512)

# Enable shadows (commented out to improve performance)
# vis.EnableShadows()

# Set the simulation time step
time_step = 1e-3

# Simulation loop
time = 0
while vis.Run():
    time += time_step

    # Set steering input for the rover after 1 second
    steering = 0
    if time >= 1:
        steering = (time - 1) * 0.2
    driver.SetSteering(steering)

    # Update rover dynamics
    rover.Update()

    # Render the scene
    vis.BeginScene()
    vis.Render()
    vis.EndScene()

    # Advance simulation by one time step
    system.DoStepDynamics(time_step)
\end{CodeBlock}

\begin{CodeBlock}{Curiosity (RBT) --- Turn 2 Ground Truth (Expert Reference)}
import os
import math
import numpy as np
import pychrono as chrono
import pychrono.robot as robot
from pychrono import irrlicht as chronoirr

# Create Chrono system
system = chrono.ChSystemNSC()
system.SetCollisionSystemType(
    chrono.ChCollisionSystem.Type_BULLET)
system.SetGravitationalAcceleration(chrono.ChVector3d(0, 0, -9.81))
chrono.ChCollisionModel.SetDefaultSuggestedEnvelope(
    0.0025)
chrono.ChCollisionModel.SetDefaultSuggestedMargin(0.0025)

# Create ground body with contact material and add it to the system
ground_mat = chrono.ChContactMaterialNSC()
ground = chrono.ChBodyEasyBox(20, 20, 1, 1000, True, True, ground_mat)
ground.SetPos(chrono.ChVector3d(0, 0, -0.5))  # Position the ground slightly below the origin
ground.SetFixed(True)  # Fix the ground in place
ground.GetVisualShape(0).SetTexture(
    chrono.GetChronoDataFile("textures/concrete.jpg"))
system.Add(ground)

# create a long box for rover to cross
box = chrono.ChBodyEasyBox(0.25, 5, 0.25, 1000, True, True, ground_mat)
box.SetPos(chrono.ChVector3d(0, 0, 0.0))
box.SetFixed(True)
box.GetVisualShape(0).SetTexture(
    chrono.GetChronoDataFile("textures/blue.png"))
system.Add(box)

# Create Curiosity rover and add it to the system
rover = robot.Curiosity(system)

# Create driver for rover
driver = robot.CuriosityDCMotorControl()
rover.SetDriver(driver)

# Initialize rover position and orientation
init_pos = chrono.ChVector3d(-5, 0.0, 0)
init_rot = chrono.ChQuaterniond(1, 0, 0, 0)
rover.Initialize(chrono.ChFramed(init_pos, init_rot))

# Create the Irrlicht visualization
vis = chronoirr.ChVisualSystemIrrlicht()
vis.AttachSystem(system)
vis.SetCameraVertical(chrono.CameraVerticalDir_Z)
vis.SetWindowSize(1280, 720)
vis.SetWindowTitle("Curiosity rover - Rigid terrain")
vis.Initialize()
vis.AddLogo(chrono.GetChronoDataFile(
    "logo_pychrono_alpha.png"))
vis.AddSkyBox()
vis.AddCamera(chrono.ChVector3d(0, 3, 3), chrono.ChVector3d(0, 0, 0))
vis.AddTypicalLights()
vis.AddLightWithShadow(chrono.ChVector3d(1.5, -2.5, 5.5), chrono.ChVector3d(0, 0, 0), 3, 4, 10, 40, 512)

# Enable shadows (commented out to improve performance)
# vis.EnableShadows()
# Set the simulation time step
time_step = 1e-3
# Simulation loop
time = 0
while vis.Run():
    time += time_step
    # ask rover to move forward
    driver.SetSteering(0.0)

    # Update rover dynamics
    rover.Update()

    # Render the scene
    vis.BeginScene()
    vis.Render()
    vis.EndScene()

    # Advance simulation by one time step
    system.DoStepDynamics(time_step)
\end{CodeBlock}

\begin{CodeBlock}[breakable]{Curiosity (RBT) --- Turn 3 Ground Truth (Expert Reference)}
import os
import math
import numpy as np
import pychrono as chrono
import pychrono.robot as robot
import pychrono.sensor as sens
from pychrono import irrlicht as chronoirr

# Create Chrono system
system = chrono.ChSystemNSC()
system.SetCollisionSystemType(
    chrono.ChCollisionSystem.Type_BULLET)
system.SetGravitationalAcceleration(chrono.ChVector3d(0, 0, -9.81))
chrono.ChCollisionModel.SetDefaultSuggestedEnvelope(
    0.0025)
chrono.ChCollisionModel.SetDefaultSuggestedMargin(0.0025)

# Create ground body with contact material and add it to the system
ground_mat = chrono.ChContactMaterialNSC()
ground = chrono.ChBodyEasyBox(20, 20, 1, 1000, True, True, ground_mat)
ground.SetPos(chrono.ChVector3d(0, 0, -0.5))  # Position the ground slightly below the origin
ground.SetFixed(True)  # Fix the ground in place
ground.GetVisualShape(0).SetTexture(
    chrono.GetChronoDataFile("textures/concrete.jpg"))
system.Add(ground)

# create a long box for rover to cross
box = chrono.ChBodyEasyBox(0.25, 5, 0.25, 1000, True, True, ground_mat)
box.SetPos(chrono.ChVector3d(0, 0, 0.0))
box.SetFixed(True)
box.GetVisualShape(0).SetTexture(
    chrono.GetChronoDataFile("textures/blue.png"))
system.Add(box)

# Create Curiosity rover and add it to the system
rover = robot.Curiosity(system)

# Create driver for rover
driver = robot.CuriosityDCMotorControl()
rover.SetDriver(driver)

# Initialize rover position and orientation
init_pos = chrono.ChVector3d(-5, 0.0, 0)
init_rot = chrono.ChQuaterniond(1, 0, 0, 0)
rover.Initialize(chrono.ChFramed(init_pos, init_rot))

# Create sensor manager
manager = sens.ChSensorManager(system)

# Create the lidar sensor and attach it to the rover
offset_pose = chrono.ChFramed(
        chrono.ChVector3d(3.0, 0, 1), chrono.QuatFromAngleAxis(0, chrono.ChVector3d(0, 1, 0))
    )

# lidar related parameters
# Update rate in Hz
update_rate = 5.0
# Number of horizontal and vertical samples
horizontal_samples = 800
vertical_samples = 300
# Horizontal and vertical field of view (radians)
horizontal_fov = 2 * chrono.CH_PI  # 360 degrees
max_vert_angle = chrono.CH_PI / 12
min_vert_angle = -chrono.CH_PI / 6
# Lag time
lag = 0
# Collection window for the lidar
collection_time = 1. / update_rate  # typically 1/update rate
# Radius of samples to use, 1->1 sample, 2->9 samples, 3->25 samples...
sample_radius = 2
# 3mm radius (as cited by velodyne)
divergence_angle = 0.003
# Lidar return mode
return_mode = sens.LidarReturnMode_STRONGEST_RETURN

lidar = sens.ChLidarSensor(
    rover.GetChassis().GetBody(),              # Body lidar is attached to
    update_rate,            # Scanning rate in Hz
    offset_pose,            # Offset pose
    horizontal_samples,     # Number of horizontal samples
    vertical_samples,       # Number of vertical channels
    horizontal_fov,         # Horizontal field of view
    max_vert_angle,         # Maximum vertical field of view
    min_vert_angle,         # Minimum vertical field of view
    100.0,                  # Maximum lidar range
    sens.LidarBeamShape_RECTANGULAR,  # Shape of the lidar beam
    sample_radius,          # Sample radius
    divergence_angle,       # Divergence angle
    divergence_angle,       # Divergence angle (again, typically same value)
    return_mode             # Return mode for the lidar
)
lidar.SetName("Lidar Sensor")
lidar.SetLag(lag)
lidar.SetCollectionWindow(collection_time)
lidar.PushFilter(sens.ChFilterDIAccess())
lidar.PushFilter(sens.ChFilterPCfromDepth())
lidar.PushFilter(sens.ChFilterVisualizePointCloud(640, 480, 1.0, "Lidar Point Cloud"))
manager.AddSensor(lidar)

# Create the Irrlicht visualization
vis = chronoirr.ChVisualSystemIrrlicht()
vis.AttachSystem(system)
vis.SetCameraVertical(chrono.CameraVerticalDir_Z)
vis.SetWindowSize(1280, 720)
vis.SetWindowTitle("Curiosity rover - Rigid terrain")
vis.Initialize()
vis.AddLogo(
    chrono.GetChronoDataFile("logo_pychrono_alpha.png"))
vis.AddSkyBox()
vis.AddCamera(chrono.ChVector3d(0, 3, 3), chrono.ChVector3d(0, 0, 0))
vis.AddTypicalLights()
vis.AddLightWithShadow(chrono.ChVector3d(1.5, -2.5, 5.5), chrono.ChVector3d(0, 0, 0), 3, 4, 10, 40, 512)

# Enable shadows (commented out to improve performance)
# vis.EnableShadows()

# Set the simulation time step
time_step = 1e-3

# Simulation loop
time = 0
while vis.Run():
    time += time_step

    # ask rover to move forward
    driver.SetSteering(0.0)

    # update sensor manager
    manager.Update()
    
    # Update rover dynamics
    rover.Update()

    # Render the scene
    vis.BeginScene()
    vis.Render()
    vis.EndScene()

    # Advance simulation by one time step
    system.DoStepDynamics(time_step)
\end{CodeBlock}
\begin{DiffBlock}{Curiosity (RBT) --- Turn 2 vs Turn 1 (Key Changes / Patch Summary)}
    (*@\textbf{Add obstacle:}@*)
    + box = chrono.ChBodyEasyBox(0.25, 5, 0.25, 1000, True, True, ground_mat)
    + box.SetPos(chrono.ChVector3d(0, 0, 0.0))
    + box.SetFixed(True)
    + box.GetVisualShape(0).SetTexture(chrono.GetChronoDataFile("textures/blue.png"))
    + system.Add(box)
    
    (*@\textbf{Change initial pose:}@*)
    - init_pos = chrono.ChVector3d(0, 0.2, 0)
    + init_pos = chrono.ChVector3d(-5, 0.0, 0)
    
    (*@\textbf{Driver logic:}@*)
    - steering = 0
    - if time >= 1:
    -     steering = (time - 1) * 0.2
    - driver.SetSteering(steering)
    + driver.SetSteering(0.0)
    \end{DiffBlock}
    
    
    \begin{DiffBlock}{Curiosity (RBT) --- Turn 3 vs Turn 2 (Key Changes / Patch Summary)}
    (*@\textbf{Add sensor module import:}@*)
    + import pychrono.sensor as sens
    
    (*@\textbf{Create sensor manager:}@*)
    + manager = sens.ChSensorManager(system)
    
    (*@\textbf{Add LiDAR sensor on rover chassis:}@*)
    + offset_pose = chrono.ChFramed(
    +     chrono.ChVector3d(3.0, 0, 1),
    +     chrono.QuatFromAngleAxis(0, chrono.ChVector3d(0, 1, 0))
    + )
    + update_rate = 5.0
    + horizontal_samples = 800
    + vertical_samples = 300
    + horizontal_fov = 2 * chrono.CH_PI
    + max_vert_angle = chrono.CH_PI / 12
    + min_vert_angle = -chrono.CH_PI / 6
    + lag = 0
    + collection_time = 1. / update_rate
    + sample_radius = 2
    + divergence_angle = 0.003
    + return_mode = sens.LidarReturnMode_STRONGEST_RETURN
    +
    + lidar = sens.ChLidarSensor(
    +     rover.GetChassis().GetBody(),
    +     update_rate,
    +     offset_pose,
    +     horizontal_samples,
    +     vertical_samples,
    +     horizontal_fov,
    +     max_vert_angle,
    +     min_vert_angle,
    +     100.0,
    +     sens.LidarBeamShape_RECTANGULAR,
    +     sample_radius,
    +     divergence_angle,
    +     divergence_angle,
    +     return_mode
    + )
    + lidar.SetName("Lidar Sensor")
    + lidar.SetLag(lag)
    + lidar.SetCollectionWindow(collection_time)
    + lidar.PushFilter(sens.ChFilterDIAccess())
    + lidar.PushFilter(sens.ChFilterPCfromDepth())
    + lidar.PushFilter(sens.ChFilterVisualizePointCloud(640, 480, 1.0, "Lidar Point Cloud"))
    + manager.AddSensor(lidar)
    
    (*@\textbf{Update sensor manager in the loop:}@*)
    + manager.Update()
    \end{DiffBlock}

    \section{Example Claude Sonnet 4 Outputs with J-LLM Feedback on Curiosity Rover Turn 3}
    \label{app:claude_curiosity_turn3}
    \begin{CodeBlock}[breakable]{Curiosity (RBT) --- Turn 3 Output by Claude Sonnet 4 (verbatim; formatting normalized)}
      import os
      import math
      import numpy as np
      import pychrono as chrono
      import pychrono.robot as robot
      import pychrono.sensor as sens
      from pychrono import irrlicht as chronoirr
      
      # Create Chrono system
      system = chrono.ChSystemNSC()
      system.SetCollisionSystemType(
        chrono.ChCollisionSystem.Type_BULLET)
      system.SetGravitationalAcceleration(
        chrono.ChVector3d(0, 0, -9.81))
      chrono.ChCollisionModel.SetDefaultSuggestedEnvelope(
        0.0025)
      chrono.ChCollisionModel.SetDefaultSuggestedMargin(
        0.0025)
      
      # Create ground body with contact material and add it to the system
      ground_mat = chrono.ChContactMaterialNSC()
      ground = chrono.ChBodyEasyBox(20, 20, 1, 1000, True, True, ground_mat)
      ground.SetPos(chrono.ChVector3d(0, 0, -0.5))  # Position the ground slightly below the origin
      ground.SetFixed(True)  # Fix the ground in place
      ground.GetVisualShape(0).SetTexture(
        chrono.GetChronoDataFile("textures/concrete.jpg"))
      system.Add(ground)
      
      # create a long box for rover to cross
      box = chrono.ChBodyEasyBox(0.25, 5, 0.25, 1000, True, True, ground_mat)
      box.SetPos(chrono.ChVector3d(0, 0, 0.0))
      box.SetFixed(True)
      box.GetVisualShape(0).SetTexture(
        chrono.GetChronoDataFile("textures/blue.png"))
      system.Add(box)
      
      # Create Curiosity rover and add it to the system
      rover = robot.Curiosity(system)
      
      # Create driver for rover
      driver = robot.CuriosityDCMotorControl()
      rover.SetDriver(driver)
      
      # Initialize rover position and orientation
      init_pos = chrono.ChVector3d(-5, 0.0, 0)
      init_rot = chrono.ChQuaterniond(1, 0, 0, 0)
      rover.Initialize(chrono.ChFramed(init_pos, init_rot))
      
      # Create sensor manager
      manager = sens.ChSensorManager(system)
      
      # Create and configure lidar sensor
      lidar_update_rate = 10  # Hz
      lidar_horizontal_samples = 800
      lidar_vertical_samples = 300
      lidar_horizontal_fov = 2 * chrono.CH_PI  # 360 degrees
      lidar_vertical_fov_max = chrono.CH_PI / 6  # 30 degrees up
      lidar_vertical_fov_min = -chrono.CH_PI / 6  # 30 degrees down
      lidar_max_distance = 100
      lidar_sample_radius = 2
      
      # Define lidar offset pose relative to rover chassis
      lidar_offset_pose = chrono.ChFramed(
          chrono.ChVector3d(0.0, 0, 1.5), chrono.QuatFromAngleAxis(0, chrono.ChVector3d(0, 1, 0))
      )
      
      # Create lidar sensor
      lidar = sens.ChLidarSensor(
          rover.GetChassis().GetBody(),  # Attach to rover chassis
          lidar_update_rate,
          lidar_offset_pose,
          lidar_horizontal_samples,
          lidar_vertical_samples,
          lidar_horizontal_fov,
          lidar_vertical_fov_max,
          lidar_vertical_fov_min,
          lidar_max_distance,
          lidar_sample_radius
      )
      lidar.SetName("Lidar Sensor")
      lidar.SetLag(0)
      lidar.SetCollectionWindow(0)
      
      # Add noise filter
      lidar.PushFilter(sens.ChFilterLidarNoiseXYZI(0.01, 0.001, 0.001, 0.01))
      
      # Add visualization filter (optional)
      lidar.PushFilter(sens.ChFilterVisualize(
        lidar_horizontal_samples, lidar_vertical_samples, "Lidar Point Cloud"))
      
      # Add data access filter
      lidar.PushFilter(sens.ChFilterDIAccess())
      
      # Add lidar to sensor manager
      manager.AddSensor(lidar)
      
      # Create the Irrlicht visualization
      vis = chronoirr.ChVisualSystemIrrlicht()
      vis.AttachSystem(system)
      vis.SetCameraVertical(chrono.CameraVerticalDir_Z)
      vis.SetWindowSize(1280, 720)
      vis.SetWindowTitle("Curiosity rover - Rigid terrain")
      vis.Initialize()
      vis.AddLogo(chrono.GetChronoDataFile(
      "logo_pychrono_alpha.png"))
      vis.AddSkyBox()
      vis.AddCamera(chrono.ChVector3d(0, 3, 3), chrono.ChVector3d(0, 0, 0))
      vis.AddTypicalLights()
      vis.AddLightWithShadow(chrono.ChVector3d(1.5, -2.5, 5.5), chrono.ChVector3d(0, 0, 0), 3, 4, 10, 40, 512)
      
      # Enable shadows (commented out to improve performance)
      # vis.EnableShadows()
      
      # Set the simulation time step
      time_step = 1e-3
      
      # Simulation loop
      time = 0
      while vis.Run():
          time += time_step
      
          # Ask rover to move forward - set both steering and throttle
          driver.SetSteering(0.0)
          driver.SetThrottle(0.5)  # Added throttle to make rover move forward
      
          # Update rover dynamics
          rover.Update()
      
          # Update sensor manager
          manager.Update()
      
          # Render the scene
          vis.BeginScene()
          vis.Render()
          vis.EndScene()
      
          # Advance simulation by one time step
          system.DoStepDynamics(time_step)
      \end{CodeBlock}
      % J-LLM output on this Claude script
      \begin{PromptBlock}{J-LLM Output on Claude Sonnet 4 (Turn 3)}
      ### Evaluation of the Simulation Script
      
      1. **Completeness (40 points total)**
         - The provided code includes all essential components such as system initialization, body creation, and visualization. However, the lidar sensor configuration lacks some details present in the reference script, particularly in terms of parameters like `divergence_angle` and `return_mode`. 
         - **Deductions**: 10 points for missing important details in lidar configuration.
      
      2. **Correctness (30 points total)**
         - The code uses the PyChrono API correctly in most areas. However, the lidar sensor's parameters are not fully aligned with the reference script, particularly regarding the `return_mode` and `divergence_angle`, which could lead to differences in sensor behavior.
         - **Deductions**: 10 points for incorrect use of lidar parameters.
      
      3. **Code Quality (10 points total)**
         - The code is well-structured and readable, with meaningful variable names. However, there are some areas where comments could be improved to provide better context, especially around the lidar sensor setup.
         - **Deductions**: 5 points for insufficient comments.
      
      4. **Efficiency (10 points total)**
         - The code is efficient and does not contain unnecessary calculations or redundant code. However, there are minor optimization opportunities in the lidar setup that could be addressed.
         - **Deductions**: 3 points for missing optimization opportunities.
      
      5. **Error Handling and Robustness (5 points total)**
         - The code does not include any explicit error handling or checks for potential issues, which is present in the reference script. This could lead to robustness issues in a real-world scenario.
         - **Deductions**: 5 points for lack of error handling.
      
      6. **Use of Visualization Tools (5 points total)**
         - The visualization setup is mostly correct, but the commented-out shadow feature could be considered a minor issue in terms of completeness.
         - **Deductions**: 3 points for inadequate visualization setup.
      
      ### Final Score
      [[64]]
      \end{PromptBlock}